% ============================================================
%  PREAMBLE — Informe de Obra Civil · EDAR
%  Asignatura: Construcción e Instalaciones Industriales
%  Máster en Ingeniería Industrial
% ============================================================

% --- Codificación y idioma ---
\usepackage[utf8]{inputenc}
\usepackage[T1]{fontenc}
\usepackage[spanish,es-nodecimaldot]{babel}

% --- Geometría de página ---
\usepackage[
  a4paper,
  top=2.5cm, bottom=2.5cm,
  left=3cm,  right=2.5cm,
  headheight=15pt
]{geometry}

% --- Fuente ---
\usepackage{lmodern}

% --- Matemáticas ---
\usepackage{amsmath}
\usepackage{amssymb}
\usepackage{mathtools}
\usepackage{bm}           % bold math

% --- Unidades SI ---
\usepackage{siunitx}
\sisetup{
  per-mode              = symbol,
  group-separator       = {.},
  output-decimal-marker = {,}
}

% --- Gráficos y figuras ---
\usepackage{graphicx}
\usepackage{float}
\usepackage{caption}
\usepackage{subcaption}
\usepackage[dvipsnames]{xcolor}

% --- TikZ (dibujos técnicos y diagramas) ---
\usepackage{tikz}
\usepackage{pgfplots}
\pgfplotsset{compat=1.18}
\usetikzlibrary{
  patterns, arrows.meta, calc, decorations.pathmorphing,
  decorations.markings, shapes.geometric, positioning, fit,
  backgrounds, shadows,
  babel        % desactiva chars activos de babel dentro de TikZ (fix >=, <->)
}

% --- Diagrama de Gantt ---
\usepackage{pgfgantt}

% --- Tablas ---
\usepackage{booktabs}
\usepackage{multirow}
\usepackage{longtable}
\usepackage{array}
\usepackage{tabularx}
\usepackage{colortbl}

% --- Listas ---
\usepackage{enumitem}

% --- Encabezados y pies de página ---
\usepackage{fancyhdr}
\pagestyle{fancy}
\fancyhf{}
\fancyhead[L]{\small\textit{Informe de Obra Civil — EDAR}}
\fancyhead[R]{\small\textit{Máster en Ingeniería Industrial}}
\fancyfoot[C]{\thepage}
\renewcommand{\headrulewidth}{0.4pt}

% --- Índice con hiperlinks ---
\usepackage[
  colorlinks=true,
  linkcolor=NavyBlue,
  citecolor=NavyBlue,
  urlcolor=NavyBlue,
  pdftitle={Informe de Obra Civil — EDAR},
  pdfauthor={Grupo Obra Civil}
]{hyperref}
\usepackage{cleveref}

% --- Cajas de notas / advertencias ---
\usepackage{tcolorbox}
\tcbuselibrary{skins, breakable}

\newtcolorbox{notabox}[1][Nota]{
  title=#1, colback=blue!5!white, colframe=NavyBlue,
  fonttitle=\bfseries, breakable
}
\newtcolorbox{supuestobox}[1][Supuesto adoptado]{
  title=#1, colback=orange!5!white, colframe=BurntOrange,
  fonttitle=\bfseries, breakable
}
\newtcolorbox{resultadobox}[1][Resultado]{
  title=#1, colback=green!5!white, colframe=OliveGreen,
  fonttitle=\bfseries, breakable
}

% --- Bibliografía ---
\usepackage{csquotes}   % requerido por biblatex con babel
\usepackage[backend=biber, style=numeric, sorting=none]{biblatex}
\addbibresource{referencias.bib}

% --- Código fuente (por si se incluyen scripts de cálculo) ---
\usepackage{listings}
\lstset{
  basicstyle=\small\ttfamily,
  breaklines=true,
  frame=single,
  language=Python
}

% ============================================================
%  PARÁMETROS DEL GRUPO (actualizar con los DNIs reales)
% ============================================================
% Media de los DNIs del grupo → X = primera cifra, Y = segunda cifra
% DNIs del grupo (6 alumnos):
%   30248832 / 49505322 / 77936142 / 01647207 / 49275542 / 53965819
% Suma = 262.578.864  →  Media = 43.763.144  →  X = 4, Y = 3
\newcommand{\paramX}{4}
\newcommand{\paramY}{3}

% Dimensiones derivadas de X e Y
\newcommand{\alturaDesnivelMuro}{7{,}3}   % (X+3,Y) = 7,3 m de desnivel
\newcommand{\largoArqueta}{4}             % largo arqueta [m]
\newcommand{\anchoArqueta}{3}             % ancho arqueta [m]
\newcommand{\alturaArqueta}{3{,}5}        % (X+Y)/2 = 3,5 m

% ============================================================
%  COMANDOS ÚTILES
% ============================================================
\newcommand{\kNm}{\si{\kilo\newton\per\metre}}
\newcommand{\kNmm}{\si{\kilo\newton\per\metre\squared}}
\newcommand{\kNcm}{\si{\kilo\newton\per\centi\metre\squared}}
\newcommand{\Tnm}{\si{\tonne\per\metre\cubed}}
\newcommand{\MPa}{\si{\mega\pascal}}

% Ecuación numerada con descripción
\newcommand{\eqtext}[1]{\quad \text{(#1)}}

% Resaltar valor final de un cálculo
\newcommand{\resultado}[1]{\boxed{#1}}
